\zhbox{55}{66}{40}{11}{8}{\centering 位}
\zhbox{99}{66}{58}{11}{8}{\centering 名称}
\zhbox{161}{66}{39}{11}{8}{\centering 默认值}
\zhbox{207}{66}{350}{11}{8}{说明}
\zhbox{207}{84}{348}{139}{8.5}{使能输出放电通路。即使 ENPWR 为低，DISCHG 位也有效。\par
0：禁用放电通路。\\1：通过内部电阻对 $V_{OUT}$ 放电。\par
此位不影响以下条件下的输出放电：通过 I\textsuperscript{2}C 设置 $V_{OUT}$、ENPWR 关闭、EN 引脚关闭、触发输出 OVP（OVP\_MODE 使能放电通路），以及触发输入 UVLO。若 GO\_BIT=1，使能放电通路；延迟20 ms后 GO\_BIT 复位至0，禁用放电通路。大多数应用建议采用较低变化率，使输出通过内部放电电流跟随 $V_{REF}$。若因 $C_{OUT}$ 过大而无法跟随，则额外增加20 ms放电时间。}
\zhbox{207}{228}{348}{39}{8.5}{使能频率抖动。\\0：禁用频率抖动。\\1：使能频率抖动。}
\zhbox{207}{271}{348}{37}{8.5}{复位 PNG 状态位，详情见第34页 PNG 位说明。\\0：$V_{OUT}$ 回到正常电压范围时，PNG 恢复为0。\\1：$V_{OUT}$ 超出电源良好电压（$V_{PG}$）范围时，PNG 锁存为1。}
\zhbox{207}{315}{348}{10}{8}{保留。应用中不得向此位写入其他值。}
\zhbox{207}{332}{348}{101}{8.5}{使能参考电压输出变化功能。写入 VREF\_L 和 VREF\_L 寄存器，然后将 GO\_BIT 写为1，$V_{REF}$ 与 $V_{OUT}$ 即由新的 VREF 值决定。主机读取 GO\_BIT 以判断参考电压调整是否完成。GO\_BIT=1时，无论 DISCHRG 位如何设置，均使能输出放电；GO\_BIT 复位为0后仍继续放电，延迟20 ms后禁用放电通路。\par
0：不能调节 $V_{OUT}$。\\1：$V_{OUT}$ 由 VREF 寄存器决定；$V_{REF}$ 达到 VREF 位设定电平后，GO\_BIT 复位为0。此设置防止错误的输出电压调整。}
\zhbox{207}{438}{348}{78}{8.5}{使能 MP2984 功率 MOSFET。ENPWR 默认值由 ADDR 引脚设置决定，详情见第28页表1。若 ENPWR 复位为0，VREF 位复位至0011 1110 100（0.5 V）。VREF 位仅应复位一次，以便 ENPWR=0后，可通过 I\textsuperscript{2}C 将 $V_{REF}$ 调整为设定值。ENPWR=0后，输出放电持续200 ms。\par
0：禁用功率 MOSFET。\\1：使能功率 MOSFET。}

